\documentclass[pdflatex,sn-chicago]{sn-jnl}

\usepackage{graphicx}
\usepackage{booktabs}
\usepackage{amsmath}
\usepackage{amssymb}
\usepackage{array}
\usepackage{caption}
\usepackage{hyperref}
\usepackage{url}
\usepackage{xcolor} % required by \email (\authemail uses \textcolor)

\title[Auditing a Released BT Evaluator for SLP]{Auditing a Released Back-Translation Evaluator for Sign Language Production: Replay Sensitivity and Public-Recipe (Non-)Reproducibility}

% LRE single-blind review: author information IS shown. Replace the bracketed
% placeholders with real names, affiliations, email and ORCID before submission.
\author*[1]{\fnm{[GivenName]} \sur{[FamilyName]}}%\orcid{0000-0000-0000-0000}
\author[1]{\fnm{[GivenName2]} \sur{[FamilyName2]}}
\email{[corresponding.author@institution.edu]}
\affil*[1]{\orgdiv{[Department]}, \orgname{[Institution]}, \orgaddress{\street{[Street]}, \city{[City]}, \postcode{[Postcode]}, \country{[Country]}}}

\begin{document}

\abstract{Back-translation (BT) evaluators mediate sign language production rankings, but a released checkpoint may yield orderings that do not survive retraining or reference replacement. On the 641-sequence PHX-public subset of SLRTP2025, a constructed whole-donor replay probe (TN-PURE-v1) scores 23.02 sacreBLEU under the released BT evaluator versus 12.78 for the recorded test poses, a gap of $+10.24$ BLEU points (95\% CI $[+8.88,+11.62]$). This constructed stress test is not reproduced by fourteen independently trained evaluators reconstructed from the public recipe: all fourteen show negative gaps (range $[-1.85,-0.80]$; 95\% prediction interval $[-1.94,-0.46]$), and across 38 non-degenerate checkpoints from further training families (dev BLEU-4 4.5--10.9, all below the released evaluator's 13.38) every gap is strictly negative, decoded with a single canonical donor registry. On the matched 461-item confidence=1 subset of Czehmann et al.'s human back-translations, reference replacement attenuates the gap from $+10.03$ to $+0.95$ (90.5\%), a reference-frame sensitivity. The released evaluator also decodes its 7{,}060-item training pool at 78.8 BLEU (70.7\% exact match) while retaining dev/test performance (13.38/12.78); permutation, input-shuffle with output-equivariance verification, stratified exact-match, and same-signer held-out controls exclude a pose-ignorant ID/cache lookup. Because competence-matched replication was not possible, these results evidence public-pipeline non-replication for a constructed stress test --- not a causal checkpoint effect, an identified memorization mechanism, or general leaderboard-rank instability.}

\keywords{sign language production, back-translation evaluation, retrieval-augmented generation, multimedia quality assessment, evaluator non-replication}

\maketitle

\section{Introduction}

Sign language production (SLP) systems are commonly evaluated by back-translation (BT): a trained sign recognizer decodes each output pose sequence, and a text metric compares the decoding with the reference. This measures recognizer-readable content through a learned measurement model, so system rankings may reflect both the evaluated outputs and the inductive biases of a particular evaluator checkpoint. The present audit is framed as a \emph{public-pipeline reproducibility audit}: our primary research question is whether the released SLRTP2025 evaluator and its replay-probe reversal can be reproduced from the publicly documented training recipe. We do \emph{not} frame this as a test of checkpoint identity, nor as a direct test of the Saunders et al.~\cite{saunders2021continuous} rank-stability proposition: that proposition is conditioned on similar model configurations, whereas every checkpoint we can construct from the public recipe is less competent than the released evaluator (dev BLEU-4 4.5--10.9 vs.\ 13.38), so competence and checkpoint identity are confounded in our design and the non-reproduction result is scoped to the public recipe.

Prior work shows that BT scores can be modest for recorded real signing and can improve without faithful target conditioning \cite{walsh2024datadriven,hong2026conditional}, and sign-native or pose-aware metrics provide complementary evidence \cite{signbleu,silverscore,jiang2025meaningful}. Nevertheless, BT remains central to benchmarks such as SLRTP2025 \cite{slrtp2025}. An unresolved question is whether a retrieval--reference stress-test ordering persists across independently trained instances of the same BT architecture.

This question is exposed by a retrieval--reference reversal on PHX-public: under the original SLRTP2025 evaluator, source-conditioned whole-donor retrieval scores 23.02 sacreBLEU (0--100 scale), whereas the recorded test poses score 12.78. This neither establishes better signing by retrieval nor invalidates BT evaluation in general --- it asks whether the reversal persists across evaluator checkpoints or depends on properties of the original evaluator.

We address this question and report three findings. \emph{First, the reversal does not replicate across independently trained evaluators reconstructed from the public recipe}: fourteen same-architecture reconstructions yield a seed-level PURE--REC gap prediction interval of $[-1.94, -0.46]$ BLEU points, and 38 non-degenerate checkpoints from further training families (distillation, config-faithful, step-faithful, rescue) extend the negative range to $[-2.01, -0.21]$. \emph{Second, the reversal is partly a property of the reference frame}: on the matched 461-item confidence=1 subset of public human sign-to-text back-translations~\cite{czehmann2026sacred}, the gap attenuates from $+10.03$ to $+0.95$ BLEU points (90.5\%), an asymmetry that is inconsistent with simple floor-effect compression. \emph{Third, the released evaluator exhibits a readout-with-generalization signature that no constructible checkpoint reproduces}: it decodes its 7{,}060-item training pool at 78.8 BLEU (70.7\% exact-match) while retaining dev/test performance of 13.38/12.78, and a Gaussian weight-perturbation sweep preserves the reversal down to dev BLEU-4 $\approx 12.7$---above the family ceiling of 10.9---before collapsing in lockstep with competence. The three findings are developed in \$Results through four lines of evidence (conditioning controls, multi-checkpoint replications, reference replacement, and five diagnostic views), with the experimental stages serving as means to the findings rather than as a linear narrative. Evaluator-aware oracle analyses are relegated to SI~Sup.~H because their unmasked advantage does not survive reference-overlap exclusion (Fig.~\ref{fig:audit_framework}).

\paragraph{Why a constructed probe matters.}
Retrieval-augmented SLP entries draw donors from the evaluator's training pool, so the TN-PURE-v1 stress test isolates a risk that production systems face in diluted form. Methodologically, the audit contributes a template --- provenance table, checkpoint registry, evaluator-replication controls --- that other benchmarks can adopt for public-recipe reproducibility testing.

\begin{figure}[t]
\centering
\includegraphics[width=\textwidth]{generated_figures/checkpoint_sensitivity_schematic.pdf}
\caption{Overview of the multi-checkpoint evaluation audit. (a) A fixed set of reference-pose, generative, retrieval, and random-donor outputs is evaluated by the original SLRTP2025 BT checkpoint and fourteen independently trained recipe-conditioned reconstructions. (b) Evaluator-aware analysis: donors selected by maximizing likelihood under one checkpoint are evaluated by all seven primary checkpoints.}
\label{fig:audit_framework}
\end{figure}

\section{Related Work}

\paragraph{Learned evaluation in sign language production.}
Neural SLP maps spoken language or glosses to pose and often evaluates with a frozen sign recognizer \cite{camgoz2018neural,saunders2020progressive}. BT was proposed as an SLP evaluation protocol by Saunders et al.~\cite{saunders2021continuous}, who argued the choice of BT model shifts absolute scores but not relative rankings between similar configurations; our audit does not directly test that proposition (its similarity condition is unmet in our reconstruction family) but characterizes a released evaluator on which a retrieval-probe reversal makes the proposition worth re-examining with matched evaluators. SignBLEU \cite{signbleu}, SiLVERScore \cite{silverscore}, and pose-native measures offer complementary views, but BT BLEU remains central to benchmarks such as SLRTP2025 \cite{slrtp2025}.

\paragraph{Retrieval and shortcut risks.}
Retrieval-based SLP ranges from phonetic composition to retrieval-enhanced generation \cite{walsh2024datadriven,signstitching,soke,darslp}. The SLRTP2025 report describes a leading retrieval-based entry, making donor copying pertinent \cite{slrtp2025}. More generally, memorization in retrieval-augmented language modeling shows that strong task scores can reflect reproduction of training examples \cite{carlini2021extracting,kandpal2022memorization,khandelwal2020knnlm}. A shortcut observed under one learned evaluator does not show that the resulting ordering persists under another checkpoint.

\paragraph{Robustness of learned metrics and corpus audits.}
Hong and Ko\v{s}eck\'{a}~\cite{hong2026conditional} show BT-BLEU can diverge from target faithfulness; Walsh et al.~\cite{walsh2024datadriven} document a low ground-truth BT baseline; Jiang et al.~\cite{jiang2025meaningful} recommend BT likelihood over decoded BLEU. Yazdani et al.~\cite{yazdani2025critical}, BackTranslation2.0~\cite{backtranslation2}, and He et al.~\cite{he2023blind} expose further fragilities of model-based generation metrics. Two 2026 PHOENIX-2014T audits bear directly on our reversal: Czehmann et al.~\cite{czehmann2026sacred} released human back-translations and documented reference information loss; Alkain et al.~\cite{alkain2026critical} showed PHOENIX-2014T has the highest train--test text overlap among tested SLT corpora. Artiaga et al.~\cite{artiaga2025signer} showed signer-dependent evaluation overestimates SLT capability. None of these studies runs multi-checkpoint evaluator experiments or donor-pool substitution controls.

\section{Methods}
\label{sec:method}

\subsection{Evaluation design}
Let $E_c$ denote a BT evaluator checkpoint and $o_{i,n}$ the pose sequence for output condition $i$ and item $n$. Checkpoint $c$ produces $h_{c,i,n}=E_c(o_{i,n})$ and $B_{c,i}=M(\{(h_{c,i,n},y_n)\}_{n=1}^{N})$, where $y_n$ is the reference text and $M$ is the corpus-level sacreBLEU functional unless stated otherwise. The values $B_{c,i}$ induce a ranking over fixed output conditions. We call a result \emph{evaluator-pipeline-sensitive} when a substantive pairwise ordering changes across independently trained checkpoints under a fixed output set. The fixed primary set comprises recorded test poses (REC), random donor copy, source-conditioned whole-donor retrieval, a Progressive Transformer generator, and retrieval--generator compositions; oracle-gloss variants are diagnostic controls.

The main evaluation uses PHX-public, a German weather-domain text-to-pose benchmark. The released SLRTP2025 pose records contain 7,060/515/641 sequences; one official test ID is absent from the released materialization, so every result is conditioned on the complete 641-record release (SI~Sup.~A). Table~\ref{tab:provenance} gives the three-way provenance: the official PHOENIX-2014T split, the released SLRTP2025 pose materialization, and the released evaluator's actual training manifest (confirmed via \texttt{config.yaml} which references the same \texttt{train.pickle}/\texttt{dev.pickle}/\texttt{test.pickle} that ship as \texttt{train.pt}/\texttt{dev.pt}/\texttt{test.pt}). Supplementary analyses use \textbf{PHX-holdout-MSKA} (133-joint, 1,538 windows) and CSL-Daily (Mode B, 1,176 sequences; SI~Sup.~L). Table~\ref{tab:protocol_summary} summarizes the protocols.

\begin{table}[t]
\centering
\scriptsize
\caption{Three-way data provenance for PHX-public. The released evaluator's training manifest matches the released pose materialization (7{,}060/515/641), not the full official PHOENIX-2014T split. All 41 missing IDs share the reason \texttt{absent\_from\_slrtp2025\_released\_pose\_records} (SI~Sup.~A). SHA-256 hashes verify file integrity.}
\label{tab:provenance}
\begin{tabular}{lrrp{3.5cm}}
\toprule
Split source & Train & Dev & Test \\
\midrule
Official PHOENIX-2014T & 7{,}096 & 519 & 642 \\
Released SLRTP2025 pose (\texttt{.pt}) & 7{,}060 & 515 & 641 \\
Released evaluator \texttt{config.yaml} manifest & 7{,}060 & 515 & 641 \\
\midrule
Missing from official & 36 & 4 & 1 \\
\bottomrule
\end{tabular}
\end{table}

\begin{table}[t]
\centering
\scriptsize
\caption{Experimental protocol summary.}
\label{tab:protocol_summary}
\begin{tabular}{llllrr}
\toprule
Setting & Recognizer & Layout & Split size & $N$ eval & $\alpha$ \\
\midrule
PHX-public  & SLRTP2025 evaluator       & 178-joint & 641 seq      & 641    & 0.88 (tuned) \\
PHX-holdout-MSKA & mska\_slr (phoenix-2014t) & 133-joint & 1{,}538 win  & 1{,}538 & 0.50 (frozen) \\
CSL-Daily   & mska (csl-daily)           & 133-joint & 256 of 1,176 win & 256    & 0.50 (frozen) \\
\bottomrule
\end{tabular}
\end{table}

All pose sequences are evaluated at 12.5\,fps, matching the SLRTP2025 evaluator training distribution; a three-path regression (canonical, explicit 25$\to$12.5-fps, and an erroneous double-subsample path) confirms that no double subsampling occurs, and that the official Codabench starter-bundle Progressive Transformer reproduces the official 1.5942 BLEU under our pipeline (full implementation details in SI~Sup.~F). PHX-public uses $\alpha{=}0.88$ (selected on dev); PHX-holdout-MSKA and CSL-Daily use $\alpha{=}0.5$.

\subsection{Systems and retrieval controls}
\label{sec:retrieval}
For each target, source-conditioned whole-donor retrieval scans the complete training pool after Unicode NFKC normalization, whitespace normalization and lowercasing, and selects the sequence with maximum source-text token-set Jaccard similarity. Oracle-gloss retrieval instead uses target gloss and is treated as an input-leakage diagnostic; random retrieval selects a donor uniformly. Donor exclusion removes exact normalized text or gloss duplicates or donors above threshold $\tau$. Gloss-based exclusion uses target gloss and is therefore an offline diagnostic rather than a deployable inference procedure.

\paragraph{Text-nearest tie-breaking.}
\label{sec:method_tie_breaking}
Jaccard ties are resolved by minimum character-level Levenshtein distance and then by minimum SHA-256 hash of the training key. The \S\ref{sec:retrieval} algorithm also excludes donors whose NFKC-normalized text equals the query's (exact-normalized-text exclusion).

\emph{Single canonical materialization.} Every number in this paper is generated from one canonical donor registry built by the deterministic \S\ref{sec:retrieval} algorithm (exact-normalized-text exclusion applied), whose SHA-256 is \texttt{9170a53026ab3263b451a3632a9e02\-318745acabf12f0b679921477afbefa301}. The artifact regenerates this registry and all headline numbers with a single command, writes them to a machine-readable \texttt{paper\_numbers.json}, and checks the manuscript against it. An earlier pre-exclusion materialization (PURE 23.79, gap $+11.01$) survives only as a sensitivity analysis in SI~Sup.~N: on 610/641 items (95.2\%) both materializations produce identical output, and a 14-reconstruction comparison gives mean $|\Delta\text{gap}| = 0.65$ BLEU (max 1.45), so qualitative conclusions are unchanged.

\paragraph{Retrieval composition.}
\label{sec:composition}
TN-PURE-v1 uses whole-donor copy with [::2] subsampling to 12.5\,fps, preserving the donor's original duration. TN-PTCOMP-v1 (a secondary control) composes a PT scaffold with donor hands and face via $\alpha$-blending of upper-body joints ($\alpha{=}0.88$ for PHX-public; full joint-set definitions and duration/word-count matching statistics in the artifact). The separately matched seen--unseen factorial retains 605/641 targets matched on duration and word count.

\paragraph{Donor-origin contrast retrieval systems.}
\label{sec:exposure_systems}
To test whether the retrieval advantage co-occurs with donor origin, we construct frozen probe systems from the same 641 queries: \textbf{UNSEEN-PURE-v1} retrieves from the 640 other test poses (excluding the query itself, its sequence root, and exact-text matches); \textbf{CROSSFIT-PURE-A/B-v1} split queries by hash into 331/310 folds, each retrieving from the other fold; \textbf{SEEN-PURE-MATCHED-v1} selects a train-pool donor within $\pm0.1$ Jaccard of each UNSEEN donor (622/641 admit one); \textbf{SEEN-RAND640-MATCHED-v1} applies matched selection inside frozen 640-donor sub-pools. SEEN-PURE-v1 re-uses the canonical TN-PURE-v1 registry. Full construction rules (calipers, tie-breaks, signer constraints) are in the artifact; these are constructed probes, not deployable generators.

\subsection{Evaluator checkpoint family}
The primary inferential checkpoint family comprises the original frozen SLRTP2025 evaluator and six \emph{recipe-conditioned reconstruction} runs trained from scratch with the documented architecture and a common frozen train/development partition (primary seeds 101--606; eight extension seeds 707--1405 use the identical protocol). The original bundle's \texttt{validations.txt} logs 202 validation points over 2{,}828 optimizer steps ($\approx$102 epochs), with decoded-BLEU selection reaching best dev BLEU-4 13.38 at step 1{,}820; intermediate checkpoints are not published. We use two reconstruction protocols: \emph{paper-derived} (14 runs, legacy script, 300 epochs, NLL selection) and \emph{released-config} (4 runs, up to 3{,}000 epochs, decoded-BLEU early stopping). Both share the documented architecture (3L/256/512ff, 8 heads), tokenizer, manifests (7{,}060/515), $[::2]$ subsampling, Adam (1e-3, betas .9/.998, wd .001), batch size 256, and plateau $\times0.8$ LR schedule (full provenance in SI). A total of 70 training runs were completed across 13 families (full registry in SI~Sup.~D). Of these, 43 were decoded on PHX-public under the canonical \S\ref{sec:retrieval} donor registry; five of these produce degenerate output (three $\alpha{=}1.0$ distillation students with 641/641 or 483/641 empty hypotheses, and two smallest ladder fractions with BLEU${}\approx 0$) and are excluded from the non-replication count, which we therefore report over 38 non-degenerate checkpoints.

\paragraph{Readout-with-generalization signature (post-hoc operational definition).}
\label{sec:readout_def}
A checkpoint exhibits the training-pool readout signature if it has both (i) training-pool free-decode BLEU $\geq 50$ \emph{or} exact-match $\geq 50\%$, and (ii) dev BLEU-4 $\geq 8.0$. These thresholds were defined after observing the released evaluator (78.8 BLEU / 70.7\% EM, dev 13.38) vs.\ the reconstruction family (max 11.5 train / 9.9 dev, no simultaneous satisfaction). The thresholds are operational heuristics for this case study, not validated criteria.

\subsection{Metrics and uncertainty}
The recorded test poses (REC) --- the released recordings of human signing, which BT decodes and compares to text references --- are decoded by each BT recognizer; we avoid the abbreviation ``ground truth'' because the recorded poses are inputs to a learned measurement model, not a semantic gold standard. We report corpus BLEU-4 on the 0--100 sacreBLEU scale as the primary metric \cite{post2018sacrebleu}; sentence-level statistics (pass-through, slot-F1) are reported separately on the 0--1 scale. BLEU-1, chrF, ROUGE-L, BERTScore \cite{zhang2020bertscore} (\texttt{bert-base-multilingual-cased}, no baseline rescaling, F1), and WER (jiwer 3.1.0 normalized) provide complementary decoded-text controls; BT likelihood and pose-DTWp provide non-decoded controls. A non-standard BLEURT substitute used only for a secondary reference-sensitivity check is documented, with its validation against canonical BLEURT-20 and its known bias, in SI~Sup.~O.

\paragraph{Scorer equivalence.}
The vendored official scorer (SacreBLEU 1.4.2) and the paper scorer (SacreBLEU 2.5.1, \texttt{tok:13a}, \texttt{smooth:exp}) give identical corpus BLEU on all 28 evaluation cells (maximum $|\Delta|=0.0000$; byte-identical hypotheses, all n-gram orders non-zero). We report the 2.5.1/13a signature as primary.

We report two forms of uncertainty. \emph{Item-sampling uncertainty} uses paired sequence-level bootstrap (10{,}000 resamples, seed 42, percentile 95\% CIs): per item, n-gram sufficient statistics are precomputed once, and every resample re-aggregates them into a corpus BLEU. Because queries reuse pose donors (565 unique donors for 641 queries), we complement the item bootstrap with donor-cluster and query$\times$donor two-way cluster bootstraps (SI~Sup.~G). A pairwise difference is called detectable when its interval excludes zero. Significance testing of the primary estimand uses a paired approximate-randomization permutation; because the headline gap was first observed in exploratory analysis and the estimand was frozen afterwards, we report a descriptive Monte-Carlo estimate: 0 of 10{,}000 permutations exceed the observed gap ($p\leq10^{-4}$, exact 95\% binomial interval $[0,0.00037]$). \emph{Training-replicate uncertainty} is reported only for the recipe-conditioned population: a descriptive seed-level 95\% Student-$t$ prediction interval (not extrapolable to the released evaluator). The eight extension seeds are post-hoc additions and enter only that interval, never the six primary comparisons. The single primary estimand is the PURE--REC decoded-sacreBLEU gap under the released evaluator; all other contrasts are secondary and are Holm-corrected within a defined family or explicitly labelled descriptive (SI~Sup.~C).

\section{Results}
\subsection{Retrieval reverses the reference ranking under the original evaluator}
\label{sec:public-diagnostic}
Under the original SLRTP2025 evaluator, source-conditioned whole-donor retrieval scored 23.02 sacreBLEU (0--100 scale), compared with 12.78 for the recorded test poses --- a retrieval--recorded-poses gap of $+10.24$ BLEU points (95\% CI $[+8.88,+11.62]$; 10,000 sequence-level resamples). This fixed-checkpoint reversal motivates the analysis but does not indicate superior sign production.

Controls localized the observation to source-conditioned donor retrieval (Table~\ref{tab:isomorphic_conditioning}): pure donor copy exceeded the oracle-gloss composed variant (11.8), canonical PT-composed (18.86), PT baseline (1.59), and 20-seed random-donor baselines (0.90 PURE, 0.77 COMP). Oracle-gloss retrieval (which uses reference gloss and is leakage) did not exceed the recorded test poses (11.4--11.8 vs 12.78). Token-normalized teacher-forced NLL (lower is better) does not reproduce the reversal: REC 3.060 vs PURE 3.479, COMP 3.411, PT 4.074. Likewise the normalized pose-DTW diagnostic does not validate semantics (REC 0 by construction vs PT/PURE/random 0.00678--0.00830). The reversal is visible under complementary text metrics on the same hypotheses (Table~\ref{tab:metric_sensitivity}): sacreBLEU-4 $+10.24$, BLEU-1 $+20.94$, chrF $+12.93$, ROUGE-L $+17.87$, WER $-13.90$ (better), BERTScore-F1 $+0.066$. We describe the observation as a \emph{decoded-text-metric reversal} and use sacreBLEU-4 as the primary decoded metric.

\paragraph{Metric implementations.}
\label{sec:metric_protocol_alignment}
All secondary metrics use the official SLRTP2025 implementations: they reproduce the official repository values (BLEU-4 12.777, chrF 34.585, WER 85.775, ROUGE 35.196 \cite{slrtp2025}) exactly. Implementation choices matter for chrF (corpus-level sacrebleu \texttt{CHRF} defaults; a sentence-level variant yields 48.05 instead of 34.59) and WER (jiwer 3.1.0 with official normalization; raw whitespace-split yields 79.26 instead of 85.78). BLEU-4 is invariant; we fix the official implementations throughout (ablations in SI).

\begin{table}[t]
\centering
\scriptsize\setlength{\tabcolsep}{3pt}
\caption{Decoded-text metric sensitivity under the original SLRTP2025 evaluator (641 sequences). The reversal direction is consistent across decoded-text metrics; teacher-forced NLL and pose-DTWp do not reproduce it (full protocol details in SI~Sup.~F).}
\label{tab:metric_sensitivity}
\begin{tabular}{lrrrrrr}
\toprule
System & BLEU-4 & BLEU-1 & chrF & ROUGE-L & WER & BERTSc. \\
\midrule
REC-v1 (recorded test poses) & 12.78 & 34.40 & 34.59 & 35.20 & 85.77 & 0.757 \\
PT-v1         & 1.59 & 16.91 & 19.03 & 17.26 & 101.68 & 0.669 \\
TN-PTCOMP  & 18.86 & 46.35 & 41.88 & 45.17 & 77.36 & 0.793 \\
TN-PURE    & \textbf{23.02} & \textbf{55.43} & \textbf{47.82} & \textbf{53.07} & \textbf{71.87} & \textbf{0.822} \\
\midrule
Gap PURE$-$REC & $+10.24$ & $+20.94$ & $+12.93$ & $+17.87$ & $-13.90$ & $+0.066$ \\
\bottomrule
\end{tabular}
\end{table}

\paragraph{Decoding-parameter sensitivity.}
Re-decoding across beam size $\{1,3,5\}$, length penalty $\{-1,0\}$, and max length $\{30,50\}$ shows the reversal is not an artifact of decoding: the PURE--REC gap is $+9.9$ (greedy), $+10.2$ (beam 3), $+10.8$ (beam 5), with length penalty and max length negligible.

\begin{table}[t]
\centering
\small
\caption{PHX-public conditioning and construction controls (641 sequences). Columns distinguish whole-donor copy from PT-scaffold composition; all rows share the query set, donor pool, evaluator, and sacreBLEU protocol. Semantic retrievers (LaBSE, multi-SBERT) lie between the canonical retriever and oracle gloss (artifact).}
\label{tab:isomorphic_conditioning}
\begin{tabular}{lrr}
\toprule
Conditioning & Pure donor copy & PT+retrieval composed \\
\midrule
Source text (German Jaccard) & \textbf{23.02} & 18.86 \\
Oracle gloss (input leakage) & 11.4 & 11.8 \\
Random donor (seeds 0--19) & 0.90 (SD 0.24) & 0.77 (SD 0.19) \\
\bottomrule
\end{tabular}
\end{table}

Pure donor copy outperformed the composed variant for every source-conditioned retriever, and source-text Jaccard produced the highest pure-donor score (23.02, ahead of LaBSE 19.2 and multi-SBERT 16.5), while oracle-gloss retrieval did not exceed the recorded test poses in either construction (11.4--11.8 vs 12.78): the reversal does not require oracle-gloss input leakage, and pure copy --- not composition --- is its strongest form.



\subsection{Independent evaluator reconstructions do not reproduce the reversal magnitude}
\label{sec:what-drives-reversal}
TN-PURE-v1 yields an original-evaluator gap of $+10.24$ BLEU points (95\% CI $[+8.88,+11.62]$). Across fourteen reconstruction runs, all fourteen have negative PURE--REC gaps (range $[-1.85, -0.80]$). TN-PTCOMP-v1 yields an original gap of $+6.09$ $[+4.76,+7.41]$; all fourteen reconstruction estimates are negative. The seed-level 95\% prediction interval is $[-1.94,-0.46]$ (mean $-1.20$, SD $0.33$, 0/14 positive). For the six primary runs, original-minus-run difference-in-differences estimates range from $+10.89$ to $+11.99$; every 95\% interval excludes zero.

\subsubsection{Recipe-conditioned evaluator reconstructions}
\label{sec:multiseed_factorial}
All fourteen reconstructions share frozen manifests, architecture, hyperparameters, and selection rule; seeds 101--1405 vary initialization and data order (Table~\ref{tab:multiseed_factorial}). All evaluators are decoded with the canonical \S\ref{sec:retrieval} donor registry (exact-normalized-text exclusion applied).

\begin{table}[h]
\centering
\scriptsize
\caption{Canonical fixed-output evaluation on the released 641-pose PHX-public subset, decoded with the \S\ref{sec:retrieval} donor registry. PURE is TN-PURE-v1; COMP is TN-PTCOMP-v1. Gaps are relative to local REC (BLEU points). Per-cell values in the artifact.}
\label{tab:multiseed_factorial}
\begin{tabular}{lrrrrr}
\toprule
Evaluator & REC & PURE & Gap & COMP & Gap \\
\midrule
Original & 12.78 & 23.02 & $\mathbf{+10.24}$ & 18.86 & $+6.09$ \\
Seed 101 & 9.24 & 8.45 & $-0.80$ & 7.92 & $-0.65$ \\
Seed 202 & 8.22 & 6.92 & $-1.30$ & 7.54 & $-1.28$ \\
Seed 303 & 9.10 & 7.25 & $-1.85$ & 8.71 & $-0.96$ \\
Seed 404 & 9.52 & 7.71 & $-1.81$ & 6.38 & $-0.81$ \\
Seed 505 & 9.71 & 8.44 & $-1.28$ & 7.97 & $-1.08$ \\
Seed 606 & 7.55 & 6.31 & $-1.24$ & 7.93 & $-0.16$ \\
\midrule
Extension 707--1405 (range) & 8.05--9.87 & 6.86--8.97 & $-1.48$ to $-0.87$ & 6.90--8.83 & $-1.03$ to $-0.44$ \\
\bottomrule
\end{tabular}
\end{table}

Competence is assessed on the released 515-sequence development subset, never on PHX-public for checkpoint selection. Original obtains dev BLEU-4 13.38 and WER 83.37 (official jiwer implementation, 0--100 scale); under the uniform full-pool beam-3 free-decode protocol, the 14 reconstruction runs obtain dev BLEU-4 6.5--9.4 and WER 88.21--94.83. Their selected validation NLL is 2.641--2.712 versus Original teacher-forced dev NLL 3.235, showing that likelihood and decoded recognition competence are not interchangeable. On PHX-public, reconstruction REC scores span 7.55--9.87 BLEU points, likewise below Original 12.78. Consequently, evaluator competence and unavailable original-training details remain confounded with evaluator identity; the experiment identifies evaluator-pipeline sensitivity, not a causal checkpoint-identity effect.

\paragraph{Competence gate and trajectory analysis.}
A checkpoint is competence-matched only if both $|\Delta\text{dev BLEU-4}|\leq 1.0$ and $|\Delta\text{dev WER}|\leq 3.0$ (the released evaluator's dev values are 13.38 / 83.37). None of the 52 non-distillation trained checkpoints passes either half (dev BLEU-4 falls short by 3.4--8.9; dev WER exceeds the 86.37 bound by 0.08--13.3). The rescue search (batch size, dropout, weight decay, label smoothing, longer budget; Table~\ref{tab:rescue_dev}) also fails: the best variant (weight-decay 0, seed 202) yields REC 9.88 / PURE 8.29 (gap $-1.59$). The gate conclusion is insensitive to threshold: across a $\pm[0.5,2.0]$ BLEU $\times$ $\pm[1.5,4.5]$ WER grid, 0 of 52 passes.

\paragraph{What the released training log establishes.}
\label{sec:original_trajectory}
The bundle's \texttt{validations.txt} logs dev BLEU-4 reaching 13.38 at step 1{,}820 (Fig.~\ref{fig:original_trajectory}), while all constructible reconstructions plateau at 8--10 by epoch 50--100. The training dynamic that produced this slope is not explained by public artifacts. Input degradation (temporal downsampling) compresses both scores but does not eliminate the reversal (at $k{=}8$: REC 8.75, PURE 11.69, gap $+2.94$; SI~Sup.~F).

\begin{figure}[t]
\centering
\includegraphics[width=0.85\textwidth]{generated_figures/original_trajectory.pdf}
\caption{Training trajectories: released evaluator (blue, decoded-BLEU selection) vs reconstruction seed 101 (red). All six primary reconstruction seeds plateau below 8.5; the released evaluator's dev perplexity bottoms around step 750 and rises thereafter, confirming decoded-BLEU rather than likelihood selection.}
\label{fig:original_trajectory}
\end{figure}

\begin{table}[h]
\centering
\small
\caption{Competence rescue sweep: 20 runs across six hyperparameter families all fail the competence gate. Best variant: weight-decay 0, seed 202 (gate threshold in \S\ref{sec:multiseed_factorial}).}
\label{tab:rescue_dev}
\begin{tabular}{lrr}
\toprule
Variant family & dev BLEU-4 range & min $|\Delta\text{BLEU-4}|$ vs Original \\
\midrule
lr $5\times10^{-4}$ / $2\times10^{-3}$ (8 runs) & 7.9--9.7 & 3.7 \\
batch size 128 / 512 & 6.8--9.2 & 4.2 \\
dropout 0.2 & 6.9--7.5 & 5.9 \\
weight decay 0 & \textbf{9.8--9.9} & \textbf{3.5} \\
label smoothing 0.1 & 6.8--7.4 & 6.0 \\
long schedule (600 epochs, lr $5\times10^{-4}$) & 8.2--8.3 & 5.1 \\
\bottomrule
\end{tabular}
\end{table}



The six primary seeds carry precomputed paired-bootstrap intervals; the eight extension seeds reproduce the same pattern (seed-level 95\% PI $[-1.94,-0.46]$).

\paragraph{Cross-metric consistency across checkpoints.}
Table~\ref{tab:cross_metric_gap} reports the gap on the same fixed output under each of the seven primary evaluators across metrics. Under the original, the gap is positive and large for all decoded-text metrics; teacher-forced NLL is more negative for PURE than REC ($-0.419$ sign-flipped), consistent with the likelihood-vs-decoded-metric divergence \cite{jiang2025meaningful}. Under the six reconstructions the reversal disappears: all chrF and ROUGE-L gaps are negative, all WER gaps favor REC, and sacreBLEU-4 gaps are all negative (range $[-1.85, -0.80]$). The reversal is a property of decoded-text metrics under the original evaluator, not specific to sacreBLEU-4.

\paragraph{Rank stability (see SI~Sup.~H).}
On a fixed nine-system stress-test set, the released evaluator ranks TN-PURE first; five of six reconstructions rank REC first, with moderate cross-checkpoint agreement (mean $\tau_b{=}0.44$, $\rho{=}0.57$).

\begin{table}[t]
\centering
\scriptsize\setlength{\tabcolsep}{3pt}
\caption{PURE--REC gap across metrics and evaluators on PHX-public (641 sequences; canonical decoding) for the original evaluator and the \emph{six primary} reconstructions (of fourteen; the eight extension seeds show the same near-zero/negative pattern in Table~\ref{tab:multiseed_factorial}). chrF, WER and ROUGE-L use the official SLRTP2025 implementations (\S\ref{sec:metric_protocol_alignment}). WER and NLL are sign-flipped so positive always means PURE is better than REC. BERTScore-F1 uses \texttt{bert-base-multilingual-cased} (bert-score 0.3.13 package, whose module \texttt{\_\_version\_\_} string reads 0.3.12; \texttt{lang=de}, \texttt{idf=False}, no baseline rescaling; model revision logged in the artifact); per-checkpoint BERTScore cells are in the artifact. The large reversal is consistent in direction across decoded-text metrics under the original evaluator and disappears under all six primary reconstructions; residual reconstruction signs are mixed and near zero for n-gram metrics.}
\label{tab:cross_metric_gap}
\begin{tabular}{lrrrrrr}
\toprule
Evaluator & $\Delta$BLEU-4 & $\Delta$BLEU-1 & $\Delta$chrF & $\Delta$WER$^\dagger$ & $\Delta$NLL$^\dagger$ & $\Delta$BERTSc. \\
\midrule
Original & $+10.24$ & $+20.94$ & $+12.93$ & $+13.90$ & $-0.419$ & $+0.066$ \\
Seed 101 & $-0.80$ & $-1.23$ & $-1.44$ & $-2.95$ & $-0.219$ & $-0.005$ \\
Seed 202 & $-1.30$ & $-1.33$ & $-1.57$ & $-3.52$ & $-0.233$ & $-0.006$ \\
Seed 303 & $-1.85$ & $-1.99$ & $-2.26$ & $-1.53$ & $-0.240$ & $+0.001$ \\
Seed 404 & $-1.81$ & $-1.60$ & $-1.90$ & $-1.87$ & $-0.194$ & $-0.001$ \\
Seed 505 & $-1.28$ & $-1.53$ & $-1.59$ & $-3.41$ & $-0.268$ & $-0.004$ \\
Seed 606 & $-1.24$ & $-1.26$ & $-1.55$ & $-3.46$ & $-0.169$ & $-0.001$ \\
\bottomrule
\end{tabular}
\par\smallskip
\footnotesize $^\dagger$WER and NLL are sign-flipped: positive means PURE is better than REC (lower WER, lower NLL).
\end{table}

\paragraph{Reference-frame sensitivity: matched-subset analysis.}
\label{sec:human_refs}
\label{sec:floor_effect}
Two 2026 corpus audits bear directly on this reversal. Czehmann et al.~\cite{czehmann2026sacred} released manual sign-to-text back-translations of the test set (one deaf fluent L2 DGS signer; per-item confidence flags), documenting information loss, lexical errors, and translationese in the original speech-transcribed references, and showed that reference replacement already changes BLEU scores and rankings on PHOENIX-2014T. Alkain et al.~\cite{alkain2026critical} quantified train--test text overlap across sign-language translation corpora, showing that a small fraction of highly training-like test items disproportionately inflates corpus BLEU. We reproduce the overlap pattern locally: the 32 test items whose reference exactly matches a training text score REC 40.26 / PURE 78.73 (gap $+38.5$), while the remaining 609 items score REC 11.93 / PURE 22.20 (gap $+10.3$).

Because Czehmann et al.\ compared original-reference and human-reference conditions on the same confidence=1 subset (462 of 642 items, excluding low-confidence, partial, or untranslatable items), we follow the same protocol. On the matched 461-item subset (after join with our 641 released test IDs), the original-reference gap is $+10.03$ BLEU points (REC 14.94 / PURE 24.98), and the human-reference gap is $+0.95$ (REC 7.35 / PURE 8.30), a paired attenuation of $9.08$ BLEU points (90.5\%, bootstrap CI $[7.44, 10.86]$; Table~\ref{tab:ref_sensitivity}; per-item paired view in Fig.~\ref{fig:ref_frame_paired}). The residual human-reference gap is not distinguishable from zero (paired item bootstrap CI $[-0.13, 1.99]$, 10{,}000 resamples), so under the independent reference frame the reversal is not merely attenuated but statistically undetectable at the item-sampling level. Teacher-forced BT likelihood does not reproduce the reversal under either reference condition (REC NLL $3.060$ vs PURE $3.479$ on original refs; REC NLL $3.41$ vs PURE $3.71$ on human refs), consistent with the likelihood-vs-decoded-metric divergence \cite{jiang2025meaningful}. Scoring against both references jointly (multi-reference, closest-length convention) leaves the gap at $+11.50$ --- \emph{larger} than the original-reference gap. This is expected under the retrieval criterion: TN-PURE donors are selected by lexical overlap with the source sentence, and the original reference is a speech transcription of that same sentence, whereas the human back-translation is an independent sign-conditioned paraphrase. Adding the original reference to the reference set therefore injects a second lexically coupled target that the replay decode can match, while the recorded poses' decode matches neither reference well; the multi-reference gap inherits the lexical coupling rather than averaging it away. The reversal is thus specific to reference frames containing the original speech-transcribed register. The full-641 comparison ($+10.24 \to +0.86$, attenuation $91.6\%$) is reported as sensitivity analysis; it mixes reference replacement with floor effects from low-confidence items and is not the primary estimand. This is a \emph{descriptive contrast of two corpus gaps}, not a causal attribution: the attenuation cannot distinguish reference-text defects, paraphrase/translationese shift, BLEU floor effects, disruption of the lexical coupling between retrieval criterion and original reference, or evaluator training-distribution alignment. Under all six reconstruction checkpoints the gap against human references is negative ($-0.15$ to $-0.65$). The asymmetry is consistent across metrics on the matched subset: chrF attenuation $82.2\%$ ($+12.73 \to +2.27$), BLEU-1 attenuation $84.0\%$ ($+20.01 \to +3.19$). Czehmann et al.\ also showed BLEURT is less sensitive to reference replacement than BLEU; our BLEURT-substitute analysis agrees and is reported with full validation details in SI~Sup.~O. We emphasize what this analysis does not show: it quantifies how the decoded text scores shift when the reference frame is changed, not whether the replay donor motion itself conveys target-equivalent meaning. The human back-translations are sign-conditioned paraphrases of the recorded test poses, not judgments of the replay donors, so the attenuation bears on the lexical scoring frame and leaves the semantic adequacy of TN-PURE-v1 unassessed.

A natural concern is that the attenuation is a symmetric floor effect. On the matched subset, switching from original to human references, REC drops $7.59$ BLEU-4 points (50.8\%) but PURE drops $16.68$ points (66.8\%) --- PURE loses a larger fraction (ratio $1.31\times$). The asymmetry is even larger under BLEU-1 and chrF, so simple symmetric compression is insufficient, but the asymmetry does not discriminate between reference-text defects and training-distribution alignment.

\begin{table}[h]
\centering
\scriptsize
\caption{Matched-subset reference sensitivity (corpus sacreBLEU, BLEU points) under the released evaluator. Primary analysis: confidence=1 subset (461 items) with original and human references computed on the \emph{same} items. Sensitivity analysis: full 641 items. Human references: Czehmann et al.~\cite{czehmann2026sacred} back-translations (CC BY-NC-SA 4.0). Confidence breakdown: 461 confidence=1.0, 121 confidence=0.5, 59 confidence=0.0.}
\label{tab:ref_sensitivity}
\begin{tabular}{lrrrr}
\toprule
Reference condition & $N$ & REC & PURE & Gap \\
\midrule
\multicolumn{5}{l}{\emph{Primary: matched confidence=1 subset}} \\
Original (speech-transcribed) & 461 & 14.94 & 24.98 & $+10.03$ \\
Human back-translation (conf=1) & 461 & 7.35 & 8.30 & $+0.95$ \\
Paired attenuation & \multicolumn{4}{r}{$9.08$ BLEU pts (90.5\%), CI $[7.44, 10.86]$} \\
\midrule
\multicolumn{5}{l}{\emph{Sensitivity: full 641 items}} \\
Original (speech-transcribed) & 641 & 12.78 & 23.02 & $+10.24$ $[+8.88,+11.62]$ \\
Human back-translation (full) & 641 & 5.66 & 6.52 & $+0.86$ $[+0.38,+1.91]$ \\
Paired attenuation & \multicolumn{4}{r}{$9.38$ BLEU pts (91.6\%)} \\
\midrule
Human refs, reconstructions (6 seeds) & 641 & 3.49--4.62 & 3.35--4.35 & $-0.15$ to $-0.65$ \\
\bottomrule
\end{tabular}
\end{table}

\begin{table}[h]
\centering
\scriptsize
\caption{Asymmetric floor-effect control (full-641 sensitivity analysis). PURE loses a larger fraction of its score than REC across all metrics when references are replaced, so simple symmetric compression is insufficient; the asymmetry does not discriminate between reference-text defects and training-distribution alignment.}
\label{tab:floor_effect}
\begin{tabular}{lrrrrrr}
\toprule
Metric & REC (orig) & REC (human) & REC drop & PURE (orig) & PURE (human) & PURE drop \\
\midrule
BLEU-4 & 12.78 & 5.66 & 7.11 (55.7\%) & 23.02 & 6.52 & 16.50 (71.7\%) \\
BLEU-1 & 34.40 & 29.28 & 5.12 (14.9\%) & 55.43 & 32.35 & 23.08 (41.6\%) \\
chrF   & 34.59 & 32.75 & 1.84 (5.3\%)  & 47.82 & 35.02 & 12.80 (26.8\%) \\
\midrule
Dual-ref gap (orig+human) & \multicolumn{6}{l}{BLEU-4: $+11.50$; BLEU-1: $+20.03$} \\
\bottomrule
\end{tabular}
\end{table}

\begin{figure}[t]
\centering
\includegraphics[width=\textwidth]{generated_figures/ref_frame_paired.pdf}
\caption{Per-item reference-frame sensitivity on the matched 461-item confidence=1 subset (released evaluator, canonical hypotheses). Sentence-level sacreBLEU (0--100, same tokenizer/smoothing as corpus BLEU). (a) Per-item scores under original vs human references for REC (recorded test poses) and PURE (replay probe); the diagonal marks equal scores. PURE points sit high under the original reference and collapse toward the diagonal/low under the human reference, whereas REC points move less. (b) Per-item paired gap (PURE$-$REC) under original (x) vs human (y) references; the star marks the mean (sentence-level mean gap $8.11 \to 0.09$, consistent in direction with the corpus-level attenuation $10.03 \to 0.95$). Sentence-level means need not equal the corpus BLEU; the figure is descriptive and does not separate reference-text defects from lexical-coupling disruption or evaluator alignment.}
\label{fig:ref_frame_paired}
\end{figure}

\subsection{Donor-origin contrast characterizes the reversal's co-occurrence with training-pool donors}
\label{sec:exposure_results}

If the released-evaluator reversal requires exposure to training-pool donors, then restricting retrieval to poses the released evaluator has not seen should remove or invert the advantage. Under the original SLRTP2025 evaluator, UNSEEN-PURE-v1 (retrieving from the 640 other test poses) scores 8.51 sacreBLEU versus 12.78 for REC, a gap of $-4.27$ BLEU points: the $+10.24$ advantage of SEEN-PURE-v1 does not survive retrieval-pool substitution. The donor-pool substitution estimand $(\text{SEEN}-\text{REC})-(\text{UNSEEN}-\text{REC})=+14.51$ BLEU points is large but is a \emph{donor-pool substitution effect}, not a pure exposure effect, because SEEN-PURE-v1 and UNSEEN-PURE-v1 differ simultaneously in pool size (7{,}060 vs 640), median Jaccard (0.50 vs 0.37), donor reuse (Gini 0.11 vs 0.29), and training-pool exposure. Cross-fit controls reproduce the same direction (CROSSFIT pooled $-3.8$ BLEU points vs fold-local REC). Under the six reconstructions all pool-origin contrasts lie in $[-0.15,+1.22]$ (Table~\ref{tab:exposure_estimands}).

\begin{table}[t]
\centering
\scriptsize\setlength{\tabcolsep}{3pt}
\caption{Donor-origin contrast estimands (sacreBLEU, BLEU points; gaps vs local REC). SEEN = 7,060 training poses; UNSEEN = 640 other test poses. Per-seed cells in the artifact.}
\label{tab:exposure_estimands}
\begin{tabular}{lrrrrr}
\toprule
Evaluator & SEEN-PURE & UNSEEN-PURE & SEEN$-$REC & UNSEEN$-$REC & Estimand \\
\midrule
Original & 23.02 & 8.51 & $+10.24$ & -4.27 & $\mathbf{+14.51}$ \\
Seeds (range) & 6.83--9.09 & 6.59--7.98 & $-1.85$ to $-0.80$ & $-1.94$ to $-0.27$ & $-0.15$ to $+1.22$ \\
\bottomrule
\end{tabular}
\end{table}

\paragraph{Path-dependent decomposition and balanced factorial.}
\label{sec:decomposition}
\label{sec:factorial_corrected}
We decompose the substitution estimand descriptively along telescoping paths (SI~Sup.~E). The path-based origin estimate varies from $+5.8$ to $+8.4$ BLEU across admissible orderings because size and selection interact. The better-balanced 605-query factorial (matching seen/unseen donors at high/low LaBSE similarity within predeclared tolerances) estimates the pool-origin main effect at $+2.08$ $[+1.56,+2.57]$, the similarity main effect at $+6.96$ $[+5.93,+8.08]$, and the seen$\times$similarity interaction at $+3.70$ $[+2.63,+4.79]$ (all Holm $p{=}0.0003$). Across reconstructions, the seen main effect shrinks to $+0.51$ and the interaction is null ($-0.39$). Two conclusions: Jaccard-only paths inflate the pool-origin component, and the released evaluator's seen advantage \emph{grows with donor--query similarity} while reconstruction interactions are null. The caveat: seen donors remain more similar than unseen within strata (SMD $+0.4$), so this is a train-pool vs test-pool provenance contrast, not an identified training-exposure effect.



\subsection{Competitive-system context from the public leaderboard}
\label{sec:competition_real}
Because the audit's stress-test set is constructed rather than competitive, the public SLRTP2025 leaderboard offers useful external context. The leaderboard was computed under the same released evaluator we audit, and its ``Ground Truth'' row reproduces our canonical REC numbers exactly (BLEU-4 12.78, BLEU-1 34.40, WER 85.77; Walsh et al.~\cite{slrtp2025}, Table~1). Table~\ref{tab:slrtp_leaderboard} reproduces the ph14t-test scores. The retrieval-based winner, USTC-MoE, sits roughly even with REC on the lexical BT metrics --- $0.72$ BLEU-4 below, $0.45$ BLEU-1 and $2.24$ chrF above, with worse WER ($93.49$ vs $85.77$) --- while its DTW-MJE ($0.0448$) and duration ratio ($1.438$) indicate considerable pose-level deviation from recorded signing. The co-occurrence of approximately matched lexical BT scores with markedly poorer pose fidelity admits several readings: that lexical BT under this evaluator is insensitive to pose fidelity, that template-dense retrieval is well matched to PHX-public's register, or that DTW-MJE itself undersells the production. The observation does not, by itself, establish that the released evaluator favors retrieval entries; resolving that question would require re-decoding fixed competition outputs under multiple evaluators. We note one naming distinction: the leaderboard's ``Progressive Transformer'' row (BLEU-4 5.43) is the original Saunders et al.~\cite{saunders2020progressive} model scored by the challenge organizers, whereas our PT-v1 control (BLEU-4 1.59) is the official Codabench starter-bundle checkpoint shipped with the SLRTP2025 release; the two are different checkpoints of the same architecture, and the gap between them illustrates the sensitivity of BT-BLEU to training and selection details even within a single model family.

\begin{table}[h]
\centering
\scriptsize\setlength{\tabcolsep}{3pt}
\caption{Public SLRTP2025 challenge leaderboard on ph14t-test (Walsh et al.~\cite{slrtp2025}, Table~1), computed under the same released BT evaluator we audit; the ``Ground Truth'' row matches our canonical REC exactly. USTC-MoE is roughly even with REC on lexical BT metrics (BLEU-4 $-0.72$, BLEU-1 $+0.45$, chrF $+2.24$, WER $+7.72$) despite a much larger DTW-MJE ($0.0448$) and duration ratio ($1.438$). This single-evaluator co-occurrence is descriptive and does not establish evaluator retrieval-favoring bias.}
\label{tab:slrtp_leaderboard}
\begin{tabular}{lrrrrrr}
\toprule
System & BLEU-1 & BLEU-4 & chrF & ROUGE & WER$\downarrow$ & DTW-MJE$\downarrow$ \\
\midrule
Recorded test poses (REC) & 34.40 & 12.78 & 34.59 & 35.20 & 85.77 & 0.0000 \\
Team 1 -- USTC-MoE (retrieval) & \textbf{34.85} & \textbf{12.06} & \textbf{36.83} & \textbf{36.59} & 93.49 & 0.0448 \\
Team 3 -- Hacettepe          & 30.44 & 9.59 & 29.70 & 30.64 & 88.88 & 0.0423 \\
Progressive Transformer      & 22.17 & 5.43 & 24.13 & 21.98 & 101.45 & 0.0418 \\
Team 2 -- hfut-lmc           & 16.96 & 2.05 & 25.88 & 19.77 & 147.85 & 0.0403 \\
\midrule
TN-PURE-v1 (our replay probe) & 55.43 & 23.02 & 47.82 & 53.07 & 71.87 & 0.00678 \\
\bottomrule
\end{tabular}
\end{table}

\paragraph{Scope limit and cross-evaluator re-decoding.}
The public challenge artifacts include only the Progressive Transformer baseline output (\texttt{PT\_baseline\_test.pt}, which we evaluate as PT-v1); the per-team pose outputs are not archived, so the competition systems cannot be re-decoded under our reconstruction family. Whether the roughly even lexical-BT scores of USTC-MoE and REC survive across independently trained evaluators is therefore an open empirical question that this audit cannot resolve. We note, as a benchmark-organizer responsibility, that archiving fixed per-team outputs alongside public scores would make rank-stability audits feasible for third parties.

\subsection{Local perturbation of the released evaluator at competence parity}
\label{sec:released_perturbation}
The reviewer's competence confound (released dev BLEU-4 13.38 vs.\ reconstructions $\leq 10.9$) cannot be broken by retraining. We can, however, perturb the released checkpoint \emph{locally}: starting from the released weights, we add i.i.d.\ Gaussian noise to every parameter (excluding BatchNorm running statistics, which must remain positive) at scale $\sigma\in\{0, 10^{-3}, 3\!\times\!10^{-3}, 10^{-2}, 3\!\times\!10^{-2}\}$ with two seeds each, and re-decode PHX-public REC/PURE and the dev set under the canonical donor registry (Fig.~\ref{fig:released_perturbation}). Small $\sigma$ perturbs checkpoint identity while preserving competence, so this design operates where the 14-seed reconstruction family cannot --- at and below the released competence from the released checkpoint itself.

The reversal is preserved under competence-preserving perturbation and collapses in lockstep with competence. At $\sigma\leq 3\!\times\!10^{-3}$ the dev BLEU-4 stays at $12.76$--$13.35$ (above every reconstruction's maximum of $10.9$) and the PURE--REC gap stays at $+9.50$ to $+10.35$ (within $\pm 0.7$ of the baseline $+10.24$). The collapse happens only at $\sigma\geq 10^{-2}$, where dev BLEU-4 simultaneously degrades to $7.3$--$10.2$ (gap $+3.2$ to $+4.2$) and at $\sigma=3\!\times\!10^{-2}$ both dev BLEU-4 ($\approx 0.1$) and the gap ($\approx 0$) are extinguished together. Two conclusions follow. First, the reversal is not a fragile identity-level artifact: at competence parity it is locally robust. Second, the reversal tracks competence rather than checkpoint identity: it survives perturbations that preserve competence and disappears only when the noise destroys competence. This shifts the unobserved interval: the reconstruction family places a lower bound at dev BLEU-4 $10.9$ (no reversal), and the perturbation sweep shows preservation down to dev BLEU-4 $\approx 12.7$, narrowing the unobserved transition to roughly $[10.9, 12.7]$ rather than $[10.9, 13.4]$.

\begin{figure}[t]
\centering
\includegraphics[width=\textwidth]{generated_figures/released_perturbation.pdf}
\caption{Released-checkpoint weight perturbation (experiment A1). (a) PURE--REC gap vs.\ Gaussian weight-noise scale $\sigma$ (two seeds per $\sigma>0$; baseline at $\sigma=0$). The reversal is preserved within $\pm 0.7$ of $+10.24$ for $\sigma\leq 3\!\times\!10^{-3}$, where dev BLEU-4 stays above $12.7$ (above every reconstruction's maximum of $10.9$). (b) Dev BLEU-4, REC BLEU-4, PURE BLEU-4, and (at baseline) train-pool readout BLEU-4. The reversal and competence collapse together at $\sigma\geq 10^{-2}$; $\sigma=10^{-1}$ produces an all-empty-decoding model and is omitted. The reversal is locally robust at competence parity and tracks competence, not identity.}
\label{fig:released_perturbation}
\end{figure}

\paragraph{Epoch-checkpoint decoupling search (experiment D).}
We additionally searched the four available intermediate epoch checkpoints (\texttt{rescue/seed\_202\_wd0} and \texttt{long\_schedule/seed\_202} at epochs 25 and 50, plus their best checkpoints) for a decoupling point that would break the within-family readout--gap coupling --- e.g., elevated train-pool readout but PURE--REC gap still $\leq 0$. No such point exists: all six checkpoints stay in the (low readout, negative gap) quadrant, with train-pool BLEU $5.5$--$8.9$ ($1.6$--$2.0\%$ EM) and PURE--REC gap $-0.94$ to $-1.70$ (per-checkpoint values in \texttt{results/epoch\_decouple.json}). The within-family relationship between low readout and negative gap is therefore monotone across the entire observable trajectory.

\subsection{Per-item gap decomposition: what kind of item drives the reversal}
\label{sec:per_item_decomp}
To characterize the \emph{mechanism} rather than the identity question, we regress the per-item sentence-level PURE--REC gap $y_i$ (sacreBLEU sentence BLEU-4 under the released evaluator, canonical donor registry; mean $8.4$, SD $18.7$, $n=641$) on interpretable item-level features: donor--reference token-set Jaccard, normalized character Levenshtein between donor and reference, an exact-match-train-text flag, the number of training texts sharing the reference's 60-character prefix (template density), the number of training texts with Jaccard $\geq 0.5$ to the reference (high-overlap-train), reference text length, donor reuse count, a confidence=1 flag, and a 9-level signer one-hot block (Fig.~\ref{fig:per_item_decomp}). The full model explains $R^2=0.33$ of the per-item variance. The strongest single-feature contributions, measured by leave-one-out $R^2$ drop, are high-overlap-train ($0.093$), template density ($0.051$), and donor--reference Levenshtein ($0.039$); the donor--reference Jaccard contributes $0.007$, the exact-match flag $0.001$, and the entire signer block $0.008$ --- signer identity is essentially irrelevant.

The signs of the standardized coefficients are informative. Items where the donor is lexically close to the reference (high Levenshtein similarity, $\beta^*=+6.1$; high donor--reference Jaccard, $\beta^*=+2.9$) show \emph{larger} gaps, whereas items with many high-overlap training neighbors ($\beta^*=-7.6$) or high template density ($\beta^*=-7.3$) show \emph{smaller} gaps. The reversal is therefore largest on items where the retrieval criterion (source-text Jaccard) happens to deliver a donor that matches the reference register closely, but where the recorded poses do not independently let the recognizer recover that text --- precisely the lexical-coupling mechanism predicted by the construction (the original reference is a speech transcription of the same source sentence the donor was selected by). Signer identity and text length contribute negligibly. The unexplained $67\%$ of variance is consistent with per-item pose variability and BT-decoder stochasticity that these surface features do not capture. This decomposition does not identify a hidden-pipeline detail; it localizes the reversal to a lexical-coupling property that is constructible on PHX-public's template-dense, high train--test overlap register, consistent with the corpus-level audit of Alkain et al.~\cite{alkain2026critical}.

\begin{figure}[t]
\centering
\includegraphics[width=\textwidth]{generated_figures/per_item_gap_decomposition.pdf}
\caption{Per-item gap decomposition (experiment C). (a) Per-feature $R^2$ drop when the feature is removed from the full OLS model ($R^2=0.33$, $n=641$). High-overlap-train count and template density are the strongest single contributors; the full signer block contributes $\leq 1\%$. (b) Per-item gap vs donor--reference token-set Jaccard; red marks the 32 items whose reference text exactly matches a training text, green the remaining 609; the black line is the bin mean. The gap grows with donor--reference lexical overlap, consistent with the lexical-coupling mechanism. Outcome is sentence-level BLEU-4 (sacreBLEU 2.5.1, 13a, exp smooth); sentence-level mean need not equal corpus BLEU.}
\label{fig:per_item_decomp}
\end{figure}

\subsection{What distinguishes the released evaluator: five converging diagnostic views}
\label{sec:positive_evidence}
The first finding (negative gaps across reconstructions) raises a positive question: what property of the released evaluator do the reconstructions lack? We approach this through five diagnostic views that triangulate on a single signature rather than independently confirming it (Fig.~\ref{fig:dose_response}; Table~\ref{tab:diagnostic_summary}; full version in SI~Sup.~D). The views reuse the same released checkpoint, PHX-public data, and related donor constructions; we therefore report them as descriptive triangulation, not as five independent tests.

\begin{figure}[t]
\centering
\includegraphics[width=\textwidth]{generated_figures/dose_response.pdf}
\caption{PURE--REC gap versus dev competence (greedy decode). Each point is a trained checkpoint (14 reconstruction seeds $+$ 15 distillation students); the released evaluator is the star at dev NLL 3.06, gap $+9.9$ under this figure's greedy-decode protocol (the canonical beam-3 headline gap is $+10.24$ at dev BLEU-4 13.38). The figure is \emph{exploratory}: checkpoints come from heterogeneous training processes and are not IID replicates; the 14 beam-3 checkpoints not plotted here use a different dev-metric protocol (SI~Sup.~D).}
\label{fig:dose_response}
\end{figure}

\begin{table}[h]
\centering
\scriptsize\setlength{\tabcolsep}{3pt}
\caption{Summary of the five diagnostic views. The released evaluator sits far outside the family range on all five views. Per-view family composition in SI~Sup.~D.}
\label{tab:diagnostic_summary}
\begin{tabular}{p{3.0cm}p{1.0cm}p{2.8cm}rr}
\toprule
Diagnostic view & Family $n$ & Released evaluator & Family range & Gap to family max \\
\midrule
(i) Gap--competence & 43 & $+10.24$ at dev 13.38 & $[-2.01, 0.00]$ at dev 4.5--10.9 & $+10.24$ \\
(ii) Pass-through ratio & 24 & 3.09 (sentence) & 1.05--1.31 & $+1.78$ \\
(iii) Train-pool readout BLEU & 27 & 78.8 (70.7\% EM) & 6.6--9.3 (1--2\% EM) & $+69.5$ \\
(iv) Frame-reversal loss & 1 & $-16.1$ (PURE) & $-4.3$ (wd0 rescue) & $|11.8|$ \\
(v) Competence gate & 52 & 13.38 & 4.5--10.02 (max) & $+3.36$ \\
\bottomrule
\end{tabular}
\end{table}

\paragraph{Readout \emph{with} generalization: the distinctive combination.}
\label{sec:insample}
The released evaluator decodes its full 7{,}060-item training pool at free-decode BLEU 78.8 with 70.7\% exact-match --- strong readout --- while retaining non-trivial dev/test performance (13.38/12.78). Under the identical full-pool, beam-3 free-decode protocol, the constructible family stays far below: 14 reconstruction seeds span train BLEU 7.2--9.3 (EM 1.5--2.3\%) at dev BLEU 6.5--9.4, and 12 non-degenerate distillation students span train BLEU 6.6--8.9 (EM 1.3--1.9\%) at dev BLEU 7.0--9.0 (these free-decode dev BLEU-4 values are not comparable to the training-time dev BLEU-4 reported in Table~\ref{tab:distill_alpha}, which was recorded under the original training protocol --- checkpoints were selected by dev NLL, SI~Sup.~E --- and reaches 10.9); an overfit rescue endpoint (epoch checkpoint) reaches high prefix-level training-decode BLEU but collapses on dev NLL (SI~Sup.~F). None of the 27 non-degenerate constructible checkpoints approaches the released evaluator's 78.8/70.7\%. Its donor-content pass-through is a 2.4$\times$ outlier (ratio 3.09 vs.\ family 1.05--1.31).

\paragraph{Leakage sanity checks for the train-pool readout.}
\label{sec:leakage_checks}
The 78.8/70.7\% readout is not a trivial implementation artifact. Five controls (full data in SI~Sup.~M) rule out four classes of shortcut: teacher forcing (code-path audit confirms autoregressive beam search), output--reference misalignment (label permutation collapses BLEU from 78.6 to $\sim$1.0), pose-ignorant ID/cache lookup (input-side pose permutation with verified output equivariance collapses BLEU to $<$1, and a same-signer held-out control reaches only BLEU 12.19 / EM 3.2\% on unseen poses versus 78.8 / 70.7\% on training poses), and narrow-text memorization (high exact-match across all 9 signers and on unique training texts). The checks do not identify the training mechanism but are consistent with pose-conditioned training-pool exposure. The thresholds defining the signature were set after observing the released evaluator and are operational heuristics, not validated criteria (\S\ref{sec:readout_def}).

\paragraph{Knowledge-distillation competence ladder.}
\label{sec:distillation_ladder}
To probe the unobserved competence interval above the reconstruction ceiling (dev BLEU-4 10.0), we trained 15 distillation students from the released evaluator as teacher: five distillation strengths ($\alpha\in\{0,0.25,0.5,0.75,1.0\}$) crossed with three seeds, with loss $\mathcal{L} = (1-\alpha)\,\text{CE} + \alpha\,T^2\,\text{KL}(\text{teacher} \| \text{student})$ at $T{=}2$. All 15 completed training and beam-3 PHX-public evaluation. Table~\ref{tab:distill_alpha} shows the distilled students extend the observable range upward (dev BLEU-4 up to 10.9), yet all twelve $\alpha{<}1.0$ students have negative PURE--REC gaps ($-2.01$ to $-0.77$). The three $\alpha{=}1.0$ students produce empty or degenerate outputs (seed 101: 641/641 empty; seed 202: 483/641 empty, 158 single-token ``.''; seed 303: 641/641 empty), yielding near-zero REC and PURE BLEU and an approximately zero gap. Training-pool readout (uniform full-7{,}060 protocol) stays at 6.6--8.9 BLEU / 1.3--1.9\% EM across the twelve $\alpha{<}1.0$ students versus the released evaluator's 78.8/70.7\%. Full soft-label distillation ($\alpha{=}1.0$) degenerates into empty or near-empty outputs and does not constitute a valid transfer test. Under the tested distillation procedure (soft-label KL at $T{=}2$ across the $\alpha$ grid), the training-pool readout signature therefore does not transfer.

\paragraph{Donor-registry consistency note.}
\label{sec:retrieval_note}
All checkpoints in this paper are decoded with the single canonical \S\ref{sec:retrieval} donor registry. The earlier pre-exclusion materialization (PURE 23.79, gap $+11.01$) and its per-item provenance analysis (30 of 31 differing items explained by the exclusion rule) are reported only as a sensitivity analysis in SI~Sup.~N; a 14-reconstruction sensitivity comparison shows mean $|\Delta\text{gap}| = 0.65$ BLEU (max 1.45).

\begin{table}[h]
\centering
\scriptsize\setlength{\tabcolsep}{3pt}
\caption{Per-$\alpha$ distillation ladder summary. 15/15 students evaluated with the canonical donor registry from \S\ref{sec:retrieval}; ``dev BLEU-4 (max)'' is the training-time dev BLEU-4 recorded by the original training protocol (legacy validation-selected protocol; not comparable to the uniform full-pool free-decode dev BLEU-4 reported in \S\ref{sec:insample}; checkpoints were selected by training-time dev NLL, SI~Sup.~E). $\alpha{=}0$ = hard-label only; $\alpha{=}1$ = full soft-label KL. All $\alpha{<}1.0$ yield negative gaps; $\alpha{=}1.0$ produces empty/degenerate outputs (641/641 or 483/641 empty hypotheses). Train-pool readout uses the identical full-7{,}060, beam-3, $[::2]$-subsample protocol for every student and for the released evaluator (no 1{,}000-item-prefix mixing); the three $\alpha{=}1.0$ students decode to empty or near-empty strings on both dev and test (uniform full-pool: 515/515 or 390/515 dev hypotheses empty, 641/641 or 483/641 test hypotheses empty) and score 0.00/0.0\%; their dev BLEU-4 is degenerate and not reported.}
\label{tab:distill_alpha}
\begin{tabular}{lrrrrr}
\toprule
$\alpha$ & $n$ & dev BLEU-4 (max) & gap (range) & train readout (BLEU / EM\%) & reaches 10.5? \\
\midrule
0.0 (hard-label) & 3 & 10.51 & $[-1.71, -0.80]$ & 8.45--8.87 / 1.9\% & 1/3 \\
0.25 & 3 & 10.75 & $[-2.01, -0.77]$ & 7.81--8.63 / 1.5--1.9\% & 2/3 \\
0.5 & 3 & 10.41 & $[-1.38, -0.77]$ & 7.73--8.20 / 1.7--1.9\% & 0/3 \\
0.75 & 3 & 10.93 & $[-1.21, -0.81]$ & 6.59--7.01 / 1.3--1.4\% & 2/3 \\
1.0 (full KL) & 3 & \emph{degenerate} & $[0.00, 0.00]$ & 0.00 / 0.0\% (empty) & --- \\
\midrule
Released (full 7{,}060, beam=3) & 1 & 13.38 & $+10.24$ & 78.82 / 70.7\% & --- \\
\bottomrule
\end{tabular}
\end{table}

\paragraph{Competence extrapolation (exploratory).}
\label{sec:extrapolation}
Three fits (OLS linear, OLS quadratic, GP with RBF + white kernel) over the 43 checkpoints spanning dev BLEU-4 $[4.5, 10.9]$ all predict negative gaps at 13.38, with substantial model sensitivity (OLS quadratic flips from $-13.5$ to $+1.3$ when distillation is removed; GP upper bound stays below $+2.0$ across leave-one-family-out folds). We treat this as an exploratory visual guide (full fit table in SI~Sup.~K), not as evidence that competence alone cannot explain the reversal. The unobserved interval is narrowed to roughly $[10.9, 12.7]$ by the released-checkpoint perturbation sweep (\S\ref{sec:released_perturbation}), which shows the reversal preserved down to dev BLEU-4 $\approx 12.7$; non-linear behavior inside this interval cannot be excluded.

\paragraph{Synthesis.}
Across the five diagnostic views, the released evaluator is the only checkpoint combining high training-pool readout (78.8 BLEU, 70.7\% EM) with retained in-domain performance (13.38/12.78 dev/test). No reconstruction or distilled student reproduces either the $+10.24$ gap or this training-pool readout signature.

\subsection{Checkpoint-targeted oracle stress tests (SI-diagnosed)}
\label{sec:non-bleu-audit}
\label{sec:bt_likelihood_attack}

We treat evaluator-aware donor selection as a checkpoint-targeted oracle stress test, not as a deployable SLP method, and we report its unmasked advantage as a non-robust diagnostic: the diagonal advantage does not survive reference-overlap exclusion (SI~Sup.~H). The full cross-checkpoint transfer matrix, candidate-pool sweep, and selector/evaluator-label permutation test are likewise in SI~Sup.~H because they share target references and evaluator by design. The main text states only that the unmasked advantage is not robust to reference-overlap exclusion and that the white-box search budget (candidate-pool size) controls whether any decoded-BLEU margin is detectable.


\section{Discussion}
\paragraph{Three-sentence summary.}
A retrieval--reference reversal of $+10.24$ BLEU points under a released BT evaluator does not survive scrutiny: it disappears under independent reconstructions of the same evaluator family, attenuates by 90.5\% under independent human references, and co-occurs with a training-pool readout signature that is unique to the released checkpoint within our constructible family. The non-replication is scoped to the public training recipe because all constructible checkpoints remain below the released evaluator's competence, leaving an unobserved interval that the released-checkpoint perturbation sweep narrows but does not eliminate. These three findings reframe the reversal as a property of one released pipeline rather than a property of BT evaluation in general, and they motivate a template for rank-stability auditing that other benchmarks can adopt.

The audit yields a broad negative result and a descriptive positive characterization. The negative result is carried by the fourteen same-recipe reconstruction seeds (shared architecture, tokenizer, manifest, selection rule; seeds 101--1405), whose Student-$t$ prediction interval for the PURE--REC gap is $[-1.94, -0.46]$ BLEU points; beyond this replication family, 38 non-degenerate checkpoints from heterogeneous families (distillation, ladder, config-faithful, step-faithful, rescue, confirmation; all decoded with the canonical donor registry from \S\ref{sec:retrieval}) all have strictly negative gaps (range $[-2.01,-0.21]$) and are reported as exploratory dose--response data, not as IID replicates of a single training process (SI~Sup.~D). The positive characterization is that the released evaluator carries a high training-pool readout combined with retained in-domain performance that no public artifact in our family exhibits (\S\ref{sec:positive_evidence}). Because all constructible checkpoints have lower competence than the released evaluator (family maximum dev BLEU-4 10.9 vs.\ released 13.38), this non-reproduction is observed under a competence confound --- but the released-checkpoint perturbation sweep (\S~\ref{sec:released_perturbation}) shows the reversal is preserved down to dev BLEU-4 $\approx 12.7$ at competence parity, narrowing the unobserved transition to roughly $[10.9, 12.7]$ and ruling out a fragile identity-level artifact. An exploratory OLS/GP extrapolation is reported in SI~Sup.~K for visualization only, without inferential claims. Because Saunders et al.~\cite{saunders2021continuous} recommended BT for comparing between similar model configurations, and our evaluator family is additionally competence-unmatched (all checkpoints dev BLEU-4 $<$ 10.9 vs.\ released 13.38), our non-reproduction result is scoped to the public training recipe and does not constitute a strong test of their rank-stability proposition.

On the matched 461-item confidence=1 subset, human-reference substitution descriptively attenuates the gap by $90.5\%$ ($+10.03 \to +0.95$, paired bootstrap CI for attenuation $[7.44, 10.86]$); the asymmetric PURE-vs-REC drop ratio ($1.31\times$ fractional for BLEU-4 on the matched-461 subset; $1.29\times$ on the full-641 sensitivity set) is inconsistent with a simple symmetric compression model but does not discriminate between reference-text defects and training-distribution alignment as non-exclusive candidate mechanisms. Czehmann et al.~\cite{czehmann2026sacred} already showed that reference replacement changes BLEU scores and rankings on PHOENIX-2014T; our contribution on this axis is applying the matched-subset protocol to the released SLRTP2025 evaluator and quantifying paired attenuation, not establishing reference-set sensitivity as a novel phenomenon. Because TN-PURE-v1 reuses pose donors across queries (641 queries select 565 unique donors; 61 donors serve more than one query, maximum reuse 6; Gini 0.11), item-level resampling overstates the effective sample size; we therefore additionally resample at the donor level. The headline gap is stable under all cluster-aware resamplings (SI~Sup.~G): donor-cluster (565 clusters, CI $[8.87,11.64]$), query$\times$donor two-way ($[8.79,11.74]$), signer (9 clusters, $[9.29,11.39]$; percentile coverage is unreliable at this cluster count, so this is a sensitivity reference only), broadcast-date (297 clusters, $[8.83,11.72]$), and signer$\times$show (17 clusters, $[7.89,11.91]$) are all consistent with the query-level CI $[8.88,11.62]$; the broadcast-show split has only 2 clusters and its interval is reported in SI~Sup.~G as an illustrative sensitivity check without confirmatory weight.

\paragraph{Relation to prior work.}
These findings extend prior critiques of BT validity from output faithfulness to public-recipe reproducibility: improved BT scores need not imply target conditioning \cite{hong2026conditional}, low recorded-pose BT baselines are documented \cite{walsh2024datadriven}, and BT likelihood is a more stable correlate of human judgment than decoded BLEU \cite{jiang2025meaningful}. Our reversal is a \emph{decoded-sacreBLEU} reversal; teacher-forced BT likelihood does not reproduce it (REC NLL $3.060$ vs PURE $3.479$), consistent with that divergence. Whole-donor replay is a constructed probe, whereas SOKE-style word-level retrieval can improve pose quality under a different representation \cite{soke}; and multiple text scorers on the same BT decoding are not independent confirmation.

\paragraph{Validation frequency is an underappreciated hyperparameter.}
Step-faithful vs epoch-validation shows validation frequency alone shifts the final checkpoint by 1.5--2.0 BLEU points (SI~Sup.~F).

\paragraph{Independent reproduction with greedy decoding.}
Under greedy decode the released evaluator scores REC 12.24 and PURE 22.12 (gap $+9.88$); no trained checkpoint reproduces the reversal. Qualitative conclusions are identical to beam-3 (artifact at \texttt{results/cells\_greedy/}).

\paragraph{What would be needed to identify the mechanism.}
Identifying the mechanism requires the SLRTP2025 organizers to disclose (i) the exact training command and seeds, (ii) the data pipeline beyond $[\,{:}\,{:}2]$ subsampling, (iii) any curriculum or checkpoint averaging, and (iv) intermediate checkpoints. A leave-one-donor-out retrain sweep is computationally tractable ($\sim$1{,}400 GPU-hours) but requires the original training infrastructure, which is not publicly available.

\paragraph{Actionable recommendations for benchmark organizers.}
The audit yields three transferable lessons, scoped to the constructed replay-probe stress case on PHX-public rather than to general SLP leaderboard stability.

\emph{Lesson 1: a released evaluator is a research artifact, not an infrastructure component.} The released \texttt{config.yaml} alone was insufficient for 70 same-recipe runs to reach competence parity, and the readout-with-generalization signature could not be localized to a specific training stage because intermediate checkpoints are unpublished. Benchmarks should release (i) the exact training command and seeds, (ii) data pipeline scripts beyond the manifest, (iii) a SHA-256 manifest of training tensors, (iv) intermediate checkpoints at least every 100 epochs so that trajectory-level analyses can localize when distinctive properties emerge, and (v) the fixed per-team pose outputs alongside the public scores. The current release provides only the Progressive Transformer baseline output; without archived per-team outputs, rank-stability audits across independently trained evaluators---like the one in this paper---cannot be reproduced by third parties on real competition systems.

\emph{Lesson 2: retrieval-based entries require shortcut-robustness certification.} A donor-exclusion sweep is cheap, deterministic, and exposes the retrieval advantage. Leaderboard entries that retrieve from the training pool should attach $\tau \in \{0.3, 0.5\}$ donor-exclusion results as a shortcut-robustness certificate. Two descriptive probes further separate genuine signing from pipeline artifacts: the training-pool free-decode BLEU $\div$ dev free-decode BLEU ratio (released: 6.1; reconstructions: 1.1--1.4) flags anomalous training-pool alignment, and the input-side pose permutation with fixed IDs and output-equivariance verification (SI~Sup.~M) excludes a pose-ignorant ID/cache lookup. The released evaluator passes the latter but fails the former by a wide margin.

\emph{Lesson 3: reference validity is part of evaluation validity.} The 90.5\% attenuation under Czehmann et al.~\cite{czehmann2026sacred} back-translations shows that rankings computed against a single speech-transcribed reference can misrepresent sign-language quality. Benchmarks should adopt human back-translations as a secondary reference set, follow the matched-subset protocol (confidence=1 items) when comparing original-reference and human-reference conditions, and flag systems whose gap diverges sharply between the two. The cost is one additional reference file and a paired bootstrap; the benefit is detection of an entire class of reference-coupling artifacts that single-reference BLEU cannot surface.

\section{Limitations}
\label{sec:limitations}
The audit is confined to PHX-public, a template-dense weather domain with documented reference-validity and train--test-overlap issues~\cite{czehmann2026sacred,alkain2026critical}. The fixed nine-system stress-test set comprises replay, composition, weak generation, and random controls rather than competitive SLP systems; the public SLRTP2025 leaderboard (\S\ref{sec:competition_real}) offers only single-evaluator context --- the retrieval-based winner reaches lexical-BT scores roughly even with REC despite far poorer pose fidelity --- and the per-team pose outputs are not publicly archived, so those real systems cannot be re-decoded under the reconstruction family. No causal claim about evaluator retrieval-favoring follows from the leaderboard alone, and the audit characterizes a constructed stress case rather than a general leaderboard-rank instability across evaluators.

All 67 non-holdout checkpoints are less competent than the released evaluator (family maximum dev BLEU-4 10.9 vs.\ released 13.38); the original training history is unpublished, so competence-matched replication from the public recipe was not possible. The released-checkpoint perturbation sweep (\S~\ref{sec:released_perturbation}) narrows the unobserved transition to roughly dev BLEU-4 $[10.9, 12.7]$, but non-linear behavior inside that interval cannot be excluded. Saunders et al.~\cite{saunders2021continuous} limited their rank-stability claim to similar model configurations; our result does not directly test their claim under those conditions. The specific mechanism is not identifiable from public artifacts: neither recipe-conditioned reconstruction, rescue sweeps, config-faithful retraining, nor any tested distillation procedure (soft-label KL at $T{=}2$ across the $\alpha$ grid; the $\alpha{=}1.0$ full-KL recipe degenerates into empty outputs and is not a valid transfer test) transfers the property. The donor-pool substitution decomposition is path-dependent; pool-origin contrasts are descriptive associations, not identified training-exposure effects. The human back-translations are a single-annotator derivative resource whose high-confidence subset lacks second-annotator adjudication.

The reversal is a decoded-BT-text-metric reversal, not a sign-language quality reversal. Without target-conditioned DGS signer judgment or a validated sign-native pose metric, we cannot determine whether the retrieved donor motion conveys target-equivalent meaning; we explicitly do not claim that donor motion fails to express the target, nor that the reversal indicates dataset contamination or pose memorization.

\paragraph{Use of public data.}
All experiments use publicly released datasets (PHOENIX-2014T, CSL-Daily, SLRTP2025) and Czehmann et al.'s human back-translations. Pose tensors may retain signer identity; we release no per-signer raw poses.

\section{Conclusion}
This audit reports three findings on a $+10.24$ BLEU-point retrieval--recorded-poses reversal under a released BT evaluator on PHX-public. The reversal does not replicate across 38 non-degenerate checkpoints from independent training families (14 same-recipe reconstructions plus 24 heterogeneous-family checkpoints; gap range $[-2.01,-0.21]$, all decoded with a single canonical donor registry). It attenuates by 90.5\% on a matched confidence=1 subset when references are replaced by independent human back-translations. And it co-occurs with a training-pool readout signature (78.8 BLEU / 70.7\% EM with retained dev/test performance of 13.38/12.78) that no constructible checkpoint exhibits; input-side pose permutation with output-equivariance verification rules out a pose-ignorant ID/cache lookup.

The scope is correspondingly bounded. A Gaussian weight-perturbation sweep preserves the reversal down to dev BLEU-4 $\approx 12.7$ and collapses in lockstep with competence, narrowing---but not eliminating---the unobserved interval between the family ceiling (10.9) and the released checkpoint (13.38). The training log shows validation PPL rising from 14.84 at step 546 to 25.42 at the best-BLEU step 1820, indicating that decoded-BLEU selection chose a checkpoint with higher validation NLL than the NLL minimum. Because competence-matched replication was not possible from public artifacts, the findings evidence public-pipeline non-replication for a constructed stress test, not a causal checkpoint effect, an identified memorization mechanism, or general leaderboard-rank instability.

\section{Compliance with Ethical Standards}
\label{sec:ethics}
\paragraph{Conflict of interest.} The authors declare no conflict of interest.
\paragraph{Ethical approval.} All experiments use publicly released datasets (PHOENIX-2014T, CSL-Daily, SLRTP2025) and Czehmann et al.'s human back-translations. No animal research was conducted.
\paragraph{Funding.} [FUNDING-PLACEHOLDER: list grants or state ``The authors declare that no funds, grants, or other support were received during the preparation of this manuscript.'']
\paragraph{Data availability.} All numerical values, per-checkpoint data, decoded hypotheses, and analysis scripts are in the anonymous artifact at \url{https://anonymous.4open.science/r/rcp-audit-artifact-B314/README.md}; a persistent DOI-bearing release will be deposited upon acceptance. Upstream datasets and their licences: PHOENIX-2014T (CC BY-NC-SA 4.0, \href{https://www-i6.informatik.rwth-aachen.de/~caminati/Phoenix-2014T.zip}{PHOENIX-2014T homepage}); SLRTP2025 released pose records and the released BT checkpoint (challenge terms, \href{https://github.com/sinebat/SLRTP2025}{SLRTP2025 repository}); CSL-Daily Mode B (available for academic research under the CSL-Daily licence, \href{http://home.ustc.edu.cn/~zhouh156/dataset/csl-daily/}{CSL-Daily homepage}); Czehmann et al.~\cite{czehmann2026sacred} human back-translations (CC BY-NC-SA 4.0). Pose tensors may retain signer identity; we release no per-signer raw poses, only de-identified decoded text and aggregate metrics. The \texttt{bert-base-multilingual-cased} model is used under the Apache-2.0 licence of its release.
\appendix
\section*{Supplementary Information (SI)}
\addcontentsline{toc}{section}{Supplementary Information}
The following supplementary materials, referenced throughout the main text, are available as a separate file (\texttt{supplementary.pdf}) and in the anonymous artifact. Each appendix is annotated with the question it answers.
\begin{itemize}\itemsep0pt
\item Sup.~A: Missing Test ID Sensitivity --- could the single absent test item change the headline gap? (\texttt{sec:appendix\_missing\_id})
\item Sup.~B: Per-seed Extension Cells 707--1405 --- does the prediction interval hold beyond the six primary seeds?
\item Sup.~C: Analysis Provenance --- what is the evidentiary role and multiplicity handling of each reported analysis? (\texttt{sec:provenance})
\item Sup.~D: Canonical Checkpoint Registry --- what exactly was trained, and which families enter which claim? (\texttt{tab:checkpoint\_registry})
\item Sup.~E: Donor-Pool Substitution Decomposition --- how much of the seen-pool advantage is pool-origin versus donor similarity? (\texttt{sec:appendix\_decomposition}, \texttt{tab:path\_sensitivity})
\item Sup.~F: Config-Faithful and Perturbation Details --- does the reconstruction plateau reflect a protocol mismatch, and how robust is the reversal to weight perturbation? (\texttt{sec:appendix\_config\_faithful}, \texttt{sec:appendix\_perturbation}, \texttt{tab:perturbation})
\item Sup.~G: Cluster-Aware Bootstrap --- does the gap survive resampling by signer, broadcast-show, and date? (\texttt{sec:appendix\_cluster\_bootstrap})
\item Sup.~H: Supplementary Oracle Diagnostics --- is the evaluator-aware advantage robust to reference-overlap exclusion? (\texttt{sec:appendix\_oracle})
\item Sup.~I: Template-Density, Holdout Boundary, Donor-Exclusion --- does the reversal persist on a clean template-holdout split and under donor-text exclusion? (\texttt{sec:appendix\_template\_boundary})
\item Sup.~J: Transcript-Slot Diagnostic Details --- do slot-level semantic scores reproduce the BLEU reversal? (\texttt{sec:appendix\_slot})
\item Sup.~K: Competence Extrapolation --- could a competence gap in the unobserved interval explain the reversal? (\texttt{sec:appendix\_extrapolation})
\item Sup.~L: CSL-Daily Boundary Control --- does the readout signature generalize beyond PHX-public? (\texttt{sec:csl-daily})
\item Sup.~M: Leakage Sanity Checks --- can the train-pool readout be explained by implementation-level shortcuts? (\texttt{sec:appendix\_leakage})
\item Sup.~N: Pre-Exclusion Materialization Sensitivity --- how sensitive are the headline numbers to the donor-exclusion rule? (\texttt{sec:appendix\_materialization})
\item Sup.~O: BLEURT-Substitute Validation --- does a learned metric corroborate the BLEU finding under reference replacement? (\texttt{sec:appendix\_bleurt})
\end{itemize}

\section*{Checkpoint Registry (summary)}

\begin{table}[h]
\centering
\scriptsize\setlength{\tabcolsep}{3pt}
\caption{Canonical checkpoint registry. 70 training runs across 13 families; 38 non-degenerate checkpoints enter the non-replication count, 43 have a canonical PHX-public PURE--REC gap, and the released evaluator is the audit target. Full per-checkpoint details in SI~Sup.~D.}
\label{tab:checkpoint_registry_summary}
\begin{tabular}{p{4.2cm}p{1.0cm}p{1.5cm}p{1.4cm}p{1.4cm}p{3.0cm}}
\toprule
Family & Planned & Trained & Has gap & Has dev & Enters \\
\midrule
Reconstructions (14) & 14 & 14 & 14 & 14 & multiseed, bootstrap, PI, dose--resp. \\
Rescue (20) & 20 & 20 & 1 & 20 & gate; wd0 dose--resp. \\
Other families (21) & 21 & 21 & 14 & 20 & dose--resp. \\
Distillation (15) & 15 & 15 & 15 & 15 & ladder, dose--resp. \\
\midrule
Total & 70 & 70 & 43 & 67 & \\
Released original & --- & --- & 1 & 1 & audit target \\
\bottomrule
\end{tabular}
\end{table}

\bibliography{references}

\end{document}
